% !TEX program = xelatex
\documentclass[11pt,a4paper]{ctexart}

% ==================== 颜色 ====================
\usepackage{xcolor}
\definecolor{maincolor}{HTML}{1A7A6B}       % 主色（标题、顶栏、框边框）
\definecolor{accent}{HTML}{D4A843}          % 强调色（关键词高亮）
\definecolor{boxbg}{HTML}{F0F7F5}           % 问题框背景
\definecolor{borderblue}{HTML}{2E86AB}      % 方法框边框
\definecolor{proofbg}{HTML}{FDF8F0}         % 方法框背景
\definecolor{summarybg}{HTML}{F5F0FF}       % 总结框背景

% ==================== 页面 ====================
\usepackage[top=2.2cm, bottom=2.2cm, left=2.5cm, right=2.5cm, headheight=23.12pt]{geometry}

% ==================== 字体 ====================
\setmainfont{TeX Gyre Pagella}

% ==================== 数学 ====================
\usepackage{amsmath,amssymb,amsfonts,mathtools}

% ==================== 彩色框 ====================
\usepackage{tcolorbox}
\tcbuselibrary{skins,breakable,most}

% ==================== 页眉页脚 ====================
\usepackage{fancyhdr}
\pagestyle{fancy}
\fancyhf{}
\fancyhead[L]{\color{maincolor}\sffamily\bfseries 数学讲义}
\fancyhead[R]{\color{gray}\sffamily \today}
\fancyfoot[C]{\color{gray}\sffamily\thepage}
\renewcommand{\headrulewidth}{1.5pt}
\renewcommand{\headrule}{\hbox to\headwidth{\color{maincolor}\leaders\hrule height \headrulewidth\hfill}}

% ==================== 章节标题样式 ====================
\usepackage{titlesec}
\titleformat{\section}
  {\color{maincolor}\sffamily\LARGE\bfseries}
  {\thesection}{1em}{}
  [\vspace{0.3cm}{\color{maincolor!40}\titlerule[2pt]}]

\titleformat{\subsection}
  {\color{maincolor!80}\sffamily\Large\bfseries}
  {\thesubsection}{1em}{}

% ==================== 表格 ====================
\usepackage{array,booktabs,colortbl}

% ==================== 杂项 ====================
\usepackage{enumitem}
\setlist{nosep, leftmargin=*}
\usepackage[hidelinks]{hyperref}

% ==================== 自定义命令 ====================
\newcommand{\highlight}[1]{\colorbox{accent!20}{\sffamily\bfseries #1}}

% ==================== 自定义环境 ====================

% 1. 问题框 —— 青绿边框 + 浅绿背景
\newtcolorbox[auto counter]{problembox}[1][]{
  enhanced,
  colframe=maincolor,
  colback=boxbg,
  coltitle=white,
  fonttitle=\sffamily\bfseries,
  title={\large 问题},
  attach boxed title to top left={xshift=4mm, yshift=-2mm},
  boxed title style={colback=maincolor, sharp corners},
  sharp corners,
  boxrule=1.2pt,
  left=4mm, right=4mm, top=5mm, bottom=4mm,
  before skip=1.2em, after skip=1.2em,
  #1
}

% 2. 方法/证明框 —— 蓝边框 + 暖白背景
\newtcolorbox[auto counter]{methodbox}[2][]{
  enhanced,
  colframe=borderblue,
  colback=proofbg,
  coltitle=white,
  fonttitle=\sffamily\bfseries,
  title={#2},
  attach boxed title to top left={xshift=4mm, yshift=-2mm},
  boxed title style={colback=borderblue, sharp corners},
  sharp corners,
  boxrule=1.2pt,
  left=4mm, right=4mm, top=5mm, bottom=4mm,
  before skip=1.2em, after skip=1.2em,
  #1
}

% 3. 总结框 —— 淡紫背景
\newtcolorbox{summarybox}[1][]{
  enhanced,
  colframe=maincolor!60,
  colback=summarybg,
  coltitle=white,
  fonttitle=\sffamily\bfseries,
  title={\large $\blacktriangleright$ 总结},
  attach boxed title to top left={xshift=4mm, yshift=-2mm},
  boxed title style={colback=maincolor!70, sharp corners},
  sharp corners,
  boxrule=1.2pt,
  left=4mm, right=4mm, top=5mm, bottom=4mm,
  before skip=1.5em, after skip=1.2em,
  #1
}

% ==================== 文档开始 ====================
\begin{document}

% ---- 彩色顶栏标题 ----
\begin{tcolorbox}[
  enhanced,
  colframe=maincolor,
  colback=maincolor,
  sharp corners,
  boxrule=0pt,
  height=2.8cm,
  valign=center,
  before skip=-2.2cm,
  after skip=1.5cm,
  grow to left by=2.5cm,
  grow to right by=2.5cm,
]
  \begin{center}
    {\color{white}\sffamily\bfseries\Huge 欧拉对质数无限的证明}\\[3pt]
    {\color{white!85}\sffamily\large 欧拉乘积公式 $\displaystyle \prod_{p}\frac{1}{1-p^{-1}} = \sum_{n=1}^{\infty}\frac{1}{n}$}
  \end{center}
\end{tcolorbox}

\vspace{0.5cm}

\begin{flushright}
  {\color{gray}\sffamily \today}
\end{flushright}

% ==================== 正文 ====================

\section{问题陈述}

\begin{problembox}
证明：\highlight{质数有无限多个}。欧拉在 1737 年首次给出了一种分析学证明，将数论中的离散问题与分析学中的连续级数联系起来，开创了\highlight{解析数论}的先河。
\end{problembox}

在欧拉之前，欧几里得已给出质数无限的经典反证法（假设有限个质数，连乘加一产生新质数）。欧拉的突破在于：他借助无穷级数的收敛与发散性质来完成证明，其核心恒等式后来被黎曼推广，成为解析数论的基石。

\section{预备知识}

\subsection*{事实一：调和级数发散}

\[
  \sum_{n=1}^{\infty} \frac{1}{n} = 1 + \frac{1}{2} + \frac{1}{3} + \cdots = \infty
\]

这是微积分中最经典的发散级数之一，欧拉对此深信不疑（尽管当时极限概念尚未严格建立）。

\subsection*{事实二：质数几何级数收敛}

对于任意大于 $1$ 的质数 $p$，其几何级数收敛：
\[
  1 + \frac{1}{p} + \frac{1}{p^2} + \cdots = \frac{1}{1 - \frac{1}{p}} = \frac{1}{1 - p^{-1}}
\]

这来自无穷等比级数求和公式 $a/(1-r)$，其中 $|r| = 1/p < 1$。

% ==================== 新页 ====================
\newpage

\section{欧拉乘积公式的推导}

\begin{methodbox}{核心思路}
将\highlight{所有质数}对应的几何级数相乘。根据算术基本定理（每个正整数可唯一分解为质因数之积），展开后的乘积恰好遍历所有正整数的倒数——从而与调和级数相等。
\end{methodbox}

\subsection*{第一步：写出所有质数级数之积}

将所有质数 $p$ 对应的几何级数相乘：
\[
  \prod_{p} \left( 1 + \frac{1}{p} + \frac{1}{p^2} + \cdots \right)
\]

\subsection*{第二步：利用算术基本定理}

任意正整数 $n$ 可以唯一地分解为质因数之积。当我们展开上述无穷乘积时，每个正整数的倒数 $\frac{1}{n}$ 恰好出现一次——这就像是对所有 $n$ 做质因数分解的"生成函数"。

例如：
\begin{itemize}
  \item 取每个括号的第一项 $1$（除了 $p=2$ 取 $\frac{1}{2}$ 外）得到 $\frac{1}{2}$
  \item 取 $p=2$ 的 $\frac{1}{2^2}$、其余取 $1$，得到 $\frac{1}{4}$
  \item 取 $p=2$ 的 $\frac{1}{2}$ 和 $p=3$ 的 $\frac{1}{3}$，得到 $\frac{1}{6}$
\end{itemize}

\subsection*{第三步：得到欧拉乘积公式}

因此，展开后恰好覆盖所有正整数的倒数之和——即调和级数：
\[
  \prod_{p} \frac{1}{1 - \frac{1}{p}} = \sum_{n=1}^{\infty} \frac{1}{n}
\]

这就是\highlight{欧拉乘积公式}（Euler Product Formula），它将质数（离散对象）与调和级数（连续分析对象）联系了起来。

\section{反证法：质数无限}

\begin{methodbox}{核心思路}
利用\highlight{反证法}。假设质数只有有限个，则欧拉乘积公式左边为有限项有限值的乘积，结果必然是一个\highlight{有限常数}；但右边是发散的调和级数（趋于\highlight{无穷大}）。有限常数不可能等于无穷大，矛盾。
\end{methodbox}

\subsection*{假设}

假设质数仅有有限的 $k$ 个：$p_1, p_2, \ldots, p_k$。

\subsection*{左边：有限常数}

则欧拉乘积公式的左边变成有限项乘积：
\[
  \prod_{i=1}^{k} \frac{1}{1 - \frac{1}{p_i}} = \prod_{i=1}^{k} \frac{p_i}{p_i - 1}
\]

每一项 $\frac{p_i}{p_i - 1}$ 都是有限的正有理数，有限个这样的数相乘，结果必然是一个\highlight{有限常数}。

\subsection*{右边：无穷大}

然而，等号右边是调和级数：
\[
  \sum_{n=1}^{\infty} \frac{1}{n} = \infty
\]

调和级数的发散性是确定的数学事实。

\subsection*{矛盾}

有限常数 $=$ 无穷大，这显然不可能。因此原假设（质数有限）不成立。
\[
  \boxed{\text{有限常数} \neq \infty}
\]

\subsection*{结论}

质数必然有\highlight{无限多个}。$\blacktriangleright$

% ==================== 新页 ====================
\newpage

\section{历史意义与后续发展}

欧拉乘积公式的意义远不止于此。它首次将\highlight{数论}（质数分布）与\highlight{分析学}（无穷级数）联系在一起，而这种联系后来深刻影响了整个数学。

\subsection*{巴塞尔问题}

欧拉在稍早时解决了著名的巴塞尔问题：
\[
  \sum_{n=1}^{\infty} \frac{1}{n^2} = \frac{\pi^2}{6}
\]
这同样展现了他高超的级数技巧。

\subsection*{黎曼 $\zeta$ 函数的推广}

约一百年后，黎曼将欧拉乘积公式中的指数推广到复数域 $s$，定义了\highlight{黎曼 $\zeta$ 函数}：
\[
  \zeta(s) = \sum_{n=1}^{\infty} \frac{1}{n^s} = \prod_{p} \frac{1}{1 - p^{-s}} \qquad (\Re(s) > 1)
\]

当 $s=1$ 时，$\zeta(1)$ 即为调和级数（发散）；欧拉乘积公式正是 $\zeta(s)$ 在 $s=1$ 处的极限情形。

\subsection*{黎曼猜想}

黎曼进一步将 $\zeta(s)$ 解析延拓到整个复平面（除 $s=1$ 外），并提出了著名的\highlight{黎曼猜想}——非平凡零点全部位于 $\Re(s) = 1/2$ 的临界线上。这一猜想至今悬而未决，是现代数学最重大的未解难题之一，其根源正可追溯到欧拉在 1737 年写下的这个恒等式。

\section{总结}

\begin{summarybox}

\subsection*{证明结构}

\begin{center}
\begin{tabular}{@{}>{\sffamily\bfseries}p{3cm} p{11cm}@{}}
\toprule
\textbf{步骤} & \textbf{内容} \\
\midrule
预备知识   & 调和级数发散 + 质数几何级数收敛于 $1/(1-p^{-1})$ \\
乘积构造   & 将所有质数的几何级数相乘 \\
算术基本定理 & 展开乘积 = 遍历所有正整数倒数之和 \\
核心等式   & $\displaystyle \prod_{p}\frac{1}{1-p^{-1}} = \sum_{n=1}^{\infty}\frac{1}{n}$ \\
反证假设   & 假设质数只有有限 $k$ 个 \\
左边       & 有限项乘积 $\to$ 有限常数 \\
右边       & 调和级数 $\to \infty$ \\
矛盾       & 有限常数 $\neq \infty$ \\
结论       & 质数必然有无穷多个 \\
\bottomrule
\end{tabular}
\end{center}

\vspace{1em}

\subsection*{结论}

欧拉乘积公式不仅优雅地证明了质数无限，更是\highlight{解析数论}的起源。它架起了数论与分析学之间的桥梁，为后世黎曼 $\zeta$ 函数和黎曼猜想奠定了基础。

\begin{center}
\begin{tcolorbox}[
  enhanced, colframe=accent!80, colback=accent!10,
  sharp corners, boxrule=1.5pt, width=10cm, center,
]
  {\color{maincolor}\sffamily\bfseries\Large 质数有无穷多个}
\end{tcolorbox}
\end{center}

\end{summarybox}

\end{document}
